\section*{F.1 结论补充说明}
\subsection*{F.1.1 参数可辨识性补充}
在不同采样频率下进行对比，发现当采样间隔过大时，多组参数会呈现近似等价拟合效果，导致可辨识性下降。建议至少保证每个主要周期内有足够采样点。

\subsection*{F.1.2 算法收敛性补充}
阻尼策略不仅影响收敛成功率，也影响收敛路径平滑性。经验上，当阻尼调整因子在 $[1.5,3.0]$ 范围时，能在速度与稳定之间取得较好平衡。

\subsection*{F.1.3 计算开销补充}
在同等硬件条件下，并行化可将批量试验总时间缩短约 35\%--50\%。其中收益取决于任务粒度与I/O开销比例。

\subsection*{F.1.4 工程实施补充}
建议对生产系统配置三类阈值：预警阈值、重估阈值、停机阈值。通过分级响应可避免瞬态噪声触发不必要的重训流程。

\subsection*{F.1.5 可复现性说明}
为确保结果可复现，应固定随机种子、记录参数边界、保存迭代日志以及软件版本信息。本文附录中的模板可直接用于复现实验记录。
